% =============================================================================
%  Une page A3 paysage par diagramme UML.
%  Chaque page est ensuite découpée par pdfseparate en un fichier PDF unique.
%  Compilation : tectonic uml-individual.tex
% =============================================================================
\documentclass[12pt]{article}

\usepackage[
  paperwidth=420mm,    % A3 paysage
  paperheight=297mm,
  left=14mm, right=14mm, top=12mm, bottom=12mm
]{geometry}

\usepackage[utf8]{inputenc}
\usepackage[T1]{fontenc}
\usepackage[french]{babel}
\usepackage{lmodern}
\usepackage{microtype}
\usepackage{graphicx}
\usepackage[export]{adjustbox}
\usepackage{xcolor}
\usepackage{tcolorbox}
\usepackage{amssymb}
\tcbuselibrary{skins}

\definecolor{orionblue}{HTML}{1F4E79}
\definecolor{orionorange}{HTML}{C8702A}
\definecolor{orionbg}{HTML}{F5F8FB}

\pagestyle{empty}
\setlength{\parindent}{0pt}

% ----------------------------------------------------------------------------
%  Macro page : un seul diagramme plein cadre.
%   #1 numéro     #2 titre      #3 sous-titre / explication
%   #4 nom de base du fichier (sans .pdf)
% ----------------------------------------------------------------------------
\newcommand{\umlpage}[4]{%
  % --- En-tête -------------------------------------------------------------
  \begin{tcolorbox}[
    enhanced,
    colback=orionblue, colframe=orionblue,
    boxrule=0pt, arc=3pt,
    left=12pt, right=12pt, top=8pt, bottom=8pt,
    width=\textwidth,
  ]
  {\color{white}\Huge\bfseries Diagramme \##1 --- #2 \hfill TP1 Orion}\\[1mm]
  {\color{orionbg}\large #3}
  \end{tcolorbox}

  \vspace{4mm}

  % --- Diagramme plein cadre ----------------------------------------------
  \vspace{2mm}
  \begin{center}
    \includegraphics[
      width=\textwidth,
      height=0.78\textheight,
      keepaspectratio
    ]{build/#4}
  \end{center}
  \vspace{2mm}

  \vfill

  % --- Pied --------------------------------------------------------------
  \begin{tcolorbox}[
    enhanced,
    colback=orionbg, colframe=orionblue!30,
    boxrule=0.4pt, arc=2pt,
    left=10pt, right=10pt, top=4pt, bottom=4pt,
    width=\textwidth,
  ]
  \small
  \textbf{Source} : \texttt{doc/rapport/diagrams/#4.mmd}
  \hfill
  \textbf{Document associé} : \texttt{doc/UML.md}
  \hfill
  ENEAM 2026 --- TP-ED, TP n\textsuperscript{o}1
  \end{tcolorbox}

  \clearpage
}

\begin{document}

\umlpage{10}{Cas d'utilisation}
        {Acteurs (analyste, dirigeants, DBA, ETL scheduler, OLTP) et
         13 cas d'utilisation issus directement de l'énoncé Orion ; relations
         «include» / «extend».}
        {10-uml-usecase}

\umlpage{11}{Diagramme de classes (modèle métier)}
        {Domaine Orion complet : héritage \texttt{Person} (Employee, Customer),
         compositions hiérarchiques produits (4 niveaux) et organisationnelles
         (5 niveaux), historisation prix, agrégation carte de fidélité,
         auto-référence du manager.}
        {11-uml-class}

\umlpage{12}{Diagramme de séquence --- passage de commande}
        {Scénario nominal OLTP : tarification depuis \texttt{product\_price\_history},
         remise active depuis \texttt{product\_discount}, transaction unique
         (\texttt{sales\_order} + N lignes), branche fidélité.}
        {12-uml-sequence-order}

\umlpage{13}{Diagramme de séquence --- job ETL SCD2}
        {Chargement de \texttt{dim\_customer} en SCD type 2 :
         calcul du \texttt{row\_hash}, fermeture de la version courante,
         insertion de la nouvelle version, mise à jour du watermark et
         journalisation \texttt{etl\_run\_log}.}
        {13-uml-sequence-etl}

\umlpage{14}{Diagramme d'activité --- cycle d'une commande}
        {Du panier au COMMIT : calcul du prix, application de la remise,
         décision \emph{stock disponible ?}, choix du canal, enregistrement,
         décision \emph{adhérent fidélité ?}, confirmation.}
        {14-uml-activity-order}

\umlpage{15}{Diagramme d'activité --- pipeline ETL}
        {Trois swimlanes APScheduler : 23h00 dimensions de référence (full reload),
         23h15 dimensions SCD2 (merge), 00h00 fact\_sales incrémental
         sur \texttt{order\_date} > watermark.}
        {15-uml-activity-etl}

\umlpage{16}{Diagramme d'états --- commande}
        {Cycle de vie : Draft $\to$ Validated $\to$ PaymentPending $\to$ Paid
         $\to$ Preparing $\to$ Shipped $\to$ Delivered $\to$ Closed,
         avec branche retour / remboursement et boucles de retry paiement.}
        {16-uml-state-order}

\umlpage{17}{Diagramme d'états --- contrat employé}
        {États : Draft $\to$ Active $\to$ (Suspended $|$ Renewed $|$ Expired
         $|$ Terminated). « Courant » au sens base
         = \texttt{end\_date IS NULL OR end\_date $\geq$ today}.}
        {17-uml-state-contract}

\umlpage{18}{Diagramme de composants}
        {Cinq composants déployables : data-gen, etl, pgAdmin, orion\_oltp,
         orion\_dwh ; interfaces SQL/JDBC fournies, dépendances Python,
         configuration partagée \texttt{.env}.}
        {18-uml-component}

\umlpage{19}{Diagramme de déploiement}
        {Topologie Docker : hôte Linux $\to$ Docker Engine $\to$ réseau
         \texttt{orion-net} $\to$ 5 conteneurs (postgres-oltp, postgres-dwh,
         etl, data-gen, pgadmin) $\to$ 3 volumes persistants.}
        {19-uml-deployment}

\umlpage{20}{Diagramme de packages}
        {Trois super-packages : \texttt{ops} (5 sous-domaines OLTP),
         \texttt{dw} (dimensions, facts, etl\_meta), \texttt{app}
         (data-gen, etl, analytics) avec dépendances «use» orientées,
         sans cycle.}
        {20-uml-package}

\end{document}
